\documentclass{article}

\usepackage[italian]{babel}
\usepackage[margin=2cm, footskip=5mm]{geometry}
% questi package non sono necessari in lualatex; ref https://tex.stackexchange.com/a/413046
% \usepackage[utf8]{inputenc}
% \usepackage[T1]{fontenc}
\usepackage{enumitem}
\usepackage{hyperref}
\usepackage{titlesec}
\usepackage{soulutf8}
\usepackage{contour}
\usepackage{float}
\usepackage{graphicx}
\usepackage{fancyhdr}
\usepackage{longtable}
\usepackage[table]{xcolor}
\usepackage{titling}
\usepackage{lastpage}
\usepackage{ifthen}
\usepackage{calc}
\usepackage{minted}
\usepackage{pgfgantt}
\usepackage{subfiles}

\newlength{\imgwidth}

\newcommand\scalegraphics[1]{%
    \settowidth{\imgwidth}{\includegraphics{#1}}%
    \setlength{\imgwidth}{\minof{\imgwidth}{\textwidth}}%
    \includegraphics[width=\imgwidth]{#1}%
}

% XXX definizione dei percorsi in cui cercare immagini
\graphicspath{ {./}
    {./img/}
}

% esempio di utilizzo: \appendToGraphicspath{./img/} (un comando diverso per ogni path da includere)
% N.B.: ci DEVE essere un forward slash alla fine del path, a indicare che è una cartella.
\makeatletter
\newcommand\appendToGraphicspath[1]{%
  \g@addto@macro\Ginput@path{{#1}}%
}
\makeatother

% setup della sottolineatura
\setuldepth{Flat}
\contourlength{0.8pt}

\newcommand{\uline}[1]{%
  \ul{{\phantom{#1}}}%
  \llap{\contour{white}{#1}}%
}

% setup dei link
\hypersetup{
  colorlinks=true, % set true if you want colored links
  linktoc=all,     % set to all if you want both sections and subsections linked
  linkcolor=black, % choose some color if you want links to stand out
}

% setup di header e footer
\pagestyle{fancy}

\fancyhf{}
\fancyhead[L]{\includegraphics[width=1cm]{logo.png}}
\fancyhead[R]{\thetitle}
\fancyfoot[R]{\thepage\ di~\pageref{LastPage}}

\fancypagestyle{nopage}{%
  \fancyfoot{}%
}

\setlength{\headheight}{1.2cm}

% setup forma \paragraph e \subparagraph
\titleformat{\paragraph}[hang]{\normalfont\normalsize\bfseries}{\theparagraph}{1em}{}
\titleformat{\subparagraph}[hang]{\normalfont\normalsize\bfseries}{\thesubparagraph}{1em}{}

% setup profondità indice di default
\setcounter{secnumdepth}{5}
\setcounter{tocdepth}{5}

% shortcut per i placeholder
\newcommand{\plchold}[1]{\textit{\{#1\}}} % chktex 20

% hook per lo script che genera il glossario
\newcommand{\glossario}[1]{\underline{#1}\textsubscript{g}}

% definizione dei comandi \uso e \stato
\makeatletter
\newcommand{\setUso}[1]{%
  \newcommand{\@uso}{#1}%
}
\newcommand{\uso}{\@uso}

\newcommand{\setStato}[1]{%
  \newcommand{\@stato}{#1}%
}
\newcommand{\stato}{\@stato}

\newcommand{\setVersione}[1]{%
  \newcommand{\@versione}{#1}%
}
\newcommand{\versione}{\@versione}

\newcommand{\setResponsabile}[1]{%
  \newcommand{\@responsabile}{#1}%
}
\newcommand{\responsabile}{\@responsabile}

\newcommand{\setRedattori}[1]{%
  \newcommand{\@redattori}{#1}%
}
\newcommand{\redattori}{\@redattori}

\newcommand{\setVerificatori}[1]{%
  \newcommand{\@verificatori}{#1}%
}
\newcommand{\verificatori}{\@verificatori}

\newcommand{\setDescrizione}[1]{%
  \newcommand{\@descrizione}{#1}%
}
\newcommand{\descrizione}{\@descrizione}

\newcommand{\setModifiche}[1]{%
  \newcommand{\@modifiche}{#1}%
}

\newcommand{\modifiche}{\@modifiche}

\makeatother

% setup delle description
\setlist[description,1]{font=$\bullet$\hspace{1.5mm}, labelwidth=* leftmargin=*,labelindent=12.5mm}
\setlist[description,2]{font=$\bullet$\hspace{1.5mm}, leftmargin=*,labelindent=12.5mm}

\appendToGraphicspath{../../commons/img/}

\title{Verbale interno --- 20/12/2019}

\setResponsabile{Alessandro Rizzo}
\setRedattori{Riccardo Agatea}
\setVerificatori{
  Tobia Apolloni \\ &
  Riccardo Cestaro
}
\setUso{Interno}
\setDescrizione{Verbale dell'incontro di GruppOne del 20/12/2019}
\setModifiche{%
\cellcolor{white!80!lightgray!100} & Alessandro Rizzo & 2019--12--21 & approva documento \\%
\cellcolor{white!80!lightgray!100} & Tobia Apolloni, Riccardo Cestaro & 2019--12--21 & verifica verbale \\%
\multirow{-3}{*}{-} \cellcolor{white!80!lightgray!100} & Luca Ercole & 2019--12--20 & stendi verbale%
}

\disabilitaElencoFigure{}
\disabilitaElencoTabelle{}

\begin{document}

\thispagestyle{empty}
\pagenumbering{gobble}

\begin{center}

  \includegraphics[width=8.5cm]{\commons/img/logo.png}\\
  {\Large GruppOne - progetto "Stalker"}\\
  \vspace{1.5cm}

  {\Huge \thetitle}
  \vspace{1.5cm}

  \begin{table}[H]
    \centering

    \begin{tabular}{r|l}
      \textbf{Versione}     & \versione              \\
      \textbf{Approvazione} & \responsabile          \\
      \textbf{Redazione}    & \redattori             \\
      \textbf{Verifica}     & \verificatori          \\
      \textbf{Stato}        & \stato                 \\
      \textbf{Uso}          & \uso                   \\
      \textbf{Destinato a}  & Imola Informatica      \\
                            & GruppOne               \\
                            & Prof. Tullio Vardanega \\
                            & Prof. Riccardo Cardin  \\
    \end{tabular}
  \end{table}

  \vspace{3cm}
  \textbf{Descrizione}\\
  \descrizione\\
  \vfill
  \verb|gruppone.swe@gmail.com|
\end{center}

\newpage
\thispagestyle{nopage}

\section*{Registro delle modifiche}
\label{sec:registro_delle_modifiche}

\begin{table}[H]
  \label{tab:registro_delle_modifiche}

  \centering
  \rowcolors{2}{lightgray}{white!80!lightgray!100}

  \begin{longtable}[c]{c c c c l}
    \rowcolor{darkgray!90!}\color{white}{\textbf{Versione}} & \color{white}{\textbf{Data}} & \color{white}{\textbf{Nominativo}} & \color{white}{\textbf{Ruolo}} & \color{white}{\textbf{Descrizione}} \\\endhead
    \modifiche
  \end{longtable}
\end{table}

% section registro_delle_modifiche (end)
\newpage

\thispagestyle{nopage}
\pagenumbering{roman}
\tableofcontents

\newpage

\pagenumbering{arabic}


\section{Informazioni logistiche}%
\label{sec:informazioni_logistiche}

\begin{description}
  \item [Luogo] Torre Archimede, aula 1BC45
  \item [Data] 20/12/2019
  \item [Ora] 13:00 \symbol{8594} 17:00
\end{description}

\subsection{Membri del gruppo presenti}%
\label{sub:membri_del_gruppo_presenti}

\begin{enumerate}
  \item Riccardo Agatea
  % \item Tobia Apolloni
  \item Riccardo Cestaro
  \item Alberto Cocco
  \item Luca Ercole
  \item Alberto Gobbo
  \item Alessandro Rizzo
  \item Fabio Scettro
\end{enumerate}
% sub:membri_del_gruppo_presenti (end)
% sec:informazioni_logistiche (end)

\section{Ordine del giorno}%
\label{sec:ordine_del_giorno}

\begin{itemize}
  \item Strumenti di comunicazione
  \item Divisione documenti da scrivere
  \item Definizione milestone
\end{itemize}
% section Ordine_del_giorno (end)

\section{Strumenti di comunicazione}%
\label{sec:strumenti_di_comunicazione}

Abbiamo scelto di utilizzare \href{Slack}{slack.com} come canale di comunicazione principale.

% sec:strumenti_di_comunicazione (end)
\section{Divisione documenti da scrivere}%
\label{sec:divisione_documenti_da_scrivere}

\subsection{Norme di progetto}%
\label{sub:norme_di_progetto}

Alberto Cocco finisce di scrivere i processi primari, Luca Ercole scrive i processi secondari e i processi organizzativi.

% sub:norme_di_progetto (end)
\subsection{Analisi dei requisiti}%
\label{sub:analisi_dei_requisiti}

Dopo aver individuato i casi d'uso, Riccardo Cestaro e Alberto Gobbo li scrivono, Fabio Scettro raccoglie i requisiti nella rispettiva tabella.

% sub:analisi_dei_requisiti (end)
\subsection{Piano di progetto}%
\label{sub:piano_di_progetto}

Alessandro Rizzo scrive l'introduzione e l'analisi dei rischi, Riccardo Agatea scrive la pianificazione e il preventivo. Rimandiamo il consuntivo al termine del periodo di investimento.

% sub:piano_di_progetto (end)
\subsection{Piano di qualifica}%
\label{sub:piano_di_qualifica}
Luca Ercole scrive l'introduzione e le sezioni sulla qualità di processo e di prodotto, Tobia Apolloni scrive i test di accettazione. Rimandiamo il calcolo dei risultati delle verifiche all'accettazione dei documenti.

% sub:piano_di_qualifica (end)
% sec:divisione_documenti_da_scrivere (end)
\section{Definizione milestone}%
\label{sec:definizione_milestone}

Abbiamo fissato come termine ultimo per la stesura delle \textit{Norme di Progetto} la data del 30/12/19, per l'analisi dei requisiti il 03/01/20, per il piano di progetto e il piano di qualifica il 07/01/20.

% sec:definizione_milestone (end)
\newpage
\section{Registro delle decisioni}%
\label{sec:registro_delle_decisioni}

\begin{enumerate}
  \item Scelto canale di comunicazione principale
  \item Scelto come dividere la stesura dei documenti
  \item Scelte deadline per la stesura dei documenti
\end{enumerate}

% sec:registro_delle_decisioni (end)
\end{document}
